% chapters/15/sub-cpmk.tex
\begin{subcpmkbox}
Bab ini melanjutkan \CPMK{4} dari sisi yang tidak terlihat. Bab sebelumnya
membahas apa yang pemain lihat; bab ini membahas apa yang pemain dengar, dan
bagaimana bunyi itu dihubungkan dengan peristiwa permainan.

\vspace{0.5em}
\textbf{\SUBCPMK{4.2} Mengintegrasikan aset audio dan memicu interaksi
audio-visual berdasarkan peristiwa \istilah{game}} --- ranah C4, \CPMK{4},
\CPL{04}.

\vspace{0.5em}
{\setlength{\tabcolsep}{4pt}\footnotesize%
\begin{tabularx}{\textwidth}{@{}C{1.9cm} L{4.0cm} L{2.5cm} Y@{}}
\toprule
\headrow \textbf{Indikator} & \textbf{Rumusan} & \textbf{Dibahas pada} & \textbf{Bukti} \\
\midrule
\ind{a} & Menghubungkan \code{AudioSource} dengan peristiwa \istilah{gameplay} &
Bagian~\ref{sec:audioclip-audiosource} dan Bagian~\ref{sec:peristiwa-bunyi} &
Daftar peristiwa permainan beserta bunyi yang dipicunya, dan bukti setiap bunyi
benar-benar terdengar pada peristiwa yang tepat. \\
\addlinespace
\ind{b} & Membedakan penggunaan BGM dan SFX secara tepat &
Bagian~\ref{sec:peran-bunyi} dan Bagian~\ref{sec:peristiwa-bunyi} &
Penjelasan tertulis atas setiap berkas bunyi, mengapa ia diperlakukan sebagai
musik latar atau sebagai efek suara. \\
\addlinespace
\ind{c} & Menyesuaikan \istilah{timing} dan level volume agar umpan balik
terasa tepat &
Bagian~\ref{sec:waktu-volume} &
Nilai volume setiap tingkat dan penyesuaian nada, beserta alasan pemilihannya
dan hasil perbandingan sebelum dan sesudah. \\
\bottomrule
\end{tabularx}}
\end{subcpmkbox}
